% Sección X: Conclusiones y trabajo futuro

\section{Conclusiones y trabajo futuro}

MediFlow aporta una base analítica reproducible e interpretable para estudiar la demanda y la presión operativa del Hospital Universidad del Norte sobre agregados OMOP CDM v5.4, integrando reconstrucción de flujo, predicción con validación temporal estricta y simulaciones paramétricas bajo supuestos explícitos. En el corte documental del repositorio (\texttt{803c1a1}), el contraste riguroso entre fuentes confirma que el experimento puntual \texttt{lgbm\_final} reduce el WAPE diario global de 18,0\,\% a 10,5\,\% (mejora relativa archivada de 41,6\,\%) con sólido desempeño en Urgencias (8,4\,\%) y Rehabilitación laboral, mientras el artefacto cuantílico servido (\texttt{v2.0}) registra WAPE de 8,91\,\%, 11,08\,\% y 42,18\,\% por servicio con coberturas de 78,57\,\%, 64,29\,\% y 73,08\,\% en 182 días. En Hospitalización, el balance documentado con $N_0=0$ en 12 cortes y la concordancia de Ley de Little sobre población común (\num{277,52} frente a \num{277,47} pacientes medios) respaldan la consistencia censal tras el filtro $p_{99{,}9}$, en tanto adelantar el 35\,\% de altas vespertinas proyecta una reducción relativa de 32,8\,\% en horas rojas (9,32\,\% a 6,26\,\%) manteniendo el pico en 359 pacientes.

Estos hallazgos retrospectivos justifican una validación operativa supervisada, pero no demuestran beneficio clínico ni sustituyen mediciones físicas del hospital. Quedan cinco prioridades inmediatas: (1) regenerar el paquete web hacia \texttt{v2.1} sobre agregados autorizados habilitando \texttt{perfiles\_horarios} reales y Ley de Little alineada; (2) recalibrar cuantiles incorporando \texttt{carnaval\_y\_semana\_santa} y nivel local en el bundle, reconciliando además las métricas de largo plazo; (3) abrir, validar visualmente y republicar el tablero Power BI (PBIP) aclarando la brecha abril--julio 2026; (4) auditar los sellos de transición e instrumentar las tres marcas HIS junto con capacidad instalada y tiempos reales de atención; y (5) ejecutar un piloto en modo sombra y aplicar el borrado seguro de datos locales al término del certamen.
